% Faktor I, III, dan V dalam satu gambar: satu codebase dibangun sekali, lalu
% digabung dengan konfigurasi tiap rumah sakit menjadi lima release.
\begin{tikzpicture}[
  font=\scriptsize,
  execute at begin node={\hyphenpenalty=10000},
  kotak/.style={draw, rounded corners, align=center, minimum height=1.05cm,
                text width=2.2cm},
  rilis/.style={draw, rounded corners, align=center, minimum height=0.62cm,
                text width=2.7cm, fill=orange!12, inner sep=2pt},
  catatan/.style={draw, rounded corners, align=center, minimum height=0.62cm,
                  text width=1.9cm, fill=blue!8, inner sep=2pt},
  panah/.style={-{Stealth[length=2mm]}, thick}
]
  \node[kotak, fill=praditagreen!15] (kode) at (0,0)
    {\textbf{Satu \textit{codebase}}\\{\tiny\mdseries SIMRS di satu repositori}};
  \node[kotak, fill=blue!8] (build) at (2.9,0)
    {\textbf{\textit{Build}}\\{\tiny\mdseries satu paket siap jalan, versi 5.3}};

  \foreach \kota/\y in {Jakarta/1.6, Bekasi/0.8, Depok/0, Bogor/-0.8, Bandung/-1.6} {
    \node[rilis] (r\kota) at (6.4,\y)
      {\textit{Release} \kota\ {\tiny\mdseries 5.3 + konfigurasi}};
    \node[catatan] (k\kota) at (9.6,\y)
      {{\tiny konfigurasi \kota}};
    \draw[panah] (build.east) -- (r\kota.west);
    \draw[panah, dashed, gray] (k\kota.west) -- (r\kota.east);
  }

  \draw[panah] (kode) -- (build);

  \node[font=\tiny, gray, anchor=south] at (rJakarta.north)
    {\textit{release} = \textit{build} + konfigurasi};
  \node[font=\tiny, gray, align=center, anchor=south] at (kJakarta.north)
    {disimpan di lingkungan,\\bukan di kode};
  \node[font=\tiny, text=praditagreen, align=center, anchor=north]
    at ([yshift=-4pt]rBandung.south)
    {\textit{run}: lima lokasi menjalankan paket yang sama};
\end{tikzpicture}
